
\documentclass[aspectratio=169,10pt]{ctexbeamer}

\usetheme{Madrid}
\usecolortheme{seahorse}
\setbeamertemplate{navigation symbols}{}
\setbeamertemplate{footline}[frame number]

\usepackage{amsmath,amssymb,booktabs,array,multirow}
\usepackage{tikz}
\usetikzlibrary{arrows.meta,positioning,shapes.geometric,fit,calc}
\usepackage{listings}
\usepackage{xcolor}
\usepackage{hyperref}

\definecolor{codebg}{RGB}{248,248,248}
\definecolor{darkblue}{RGB}{0,65,120}
\definecolor{darkgreen}{RGB}{0,105,70}
\definecolor{warnorange}{RGB}{190,110,0}
\definecolor{softred}{RGB}{205,70,70}

\hypersetup{colorlinks=true,linkcolor=darkblue,urlcolor=darkblue}

\lstdefinestyle{pythonstyle}{
    language=Python,
    basicstyle=\ttfamily\scriptsize,
    keywordstyle=\color{darkblue}\bfseries,
    commentstyle=\color{darkgreen},
    stringstyle=\color{warnorange},
    showstringspaces=false,
    columns=fullflexible,
    backgroundcolor=\color{codebg},
    frame=single,
    rulecolor=\color{black!15},
    breaklines=true
}
\lstset{style=pythonstyle}

\newcommand{\code}[1]{\texttt{#1}}
\newcommand{\checkitem}{\textcolor{darkgreen}{\Large$\checkmark$}\;}
\newcommand{\warnitem}{\textcolor{warnorange}{\Large$\triangle$}\;}
\newcommand{\baditem}{\textcolor{softred}{\Large$\times$}\;}

\title[Week 2: Three-phase pandapower]{Week 2: 三相不平衡潮流建模与验证}
\subtitle{pandapower \texttt{runpp\_3ph()} for Microgrid Security Prediction}
\author{ECE Course - Research Tutorial}
\institute{Microgrid 3-phase Security AI Project}
\date{2026}

\begin{document}

% ----------------------------------------------------------------------
\begin{frame}
  \titlepage
  \vspace{-0.8em}
  \begin{center}
    \small 本周关键词: three-phase, unbalanced power flow, symmetrical components, VUF, proof, cross-validation
  \end{center}
\end{frame}

% ----------------------------------------------------------------------
\begin{frame}{本周学习目标}
\begin{columns}[T]
\column{0.54\textwidth}
学生完成 Week 2 后应能:
\begin{itemize}
  \item 解释微电网为什么需要三相不平衡潮流;
  \item 使用 \code{asymmetric\_load} 建模逐相负荷;
  \item 使用 \code{asymmetric\_sgen} 建模逐相 PV/DER;
  \item 使用 \code{pp.runpp\_3ph(net)} 运行三相潮流;
  \item 读取 \code{res\_bus\_3ph}, \code{res\_line\_3ph}, \code{res\_ext\_grid\_3ph};
  \item 用对称分量、功率平衡和线路电流公式验证结果。
\end{itemize}
\column{0.40\textwidth}
\begin{block}{本周边界}
本周先把三相 base-case 潮流算对、验对。Week 4 再把相同检查扩展为 N-1 violation label。
\end{block}
\begin{alertblock}{本周核心能力}
不是“会调用函数”, 而是能证明输出结果与基本物理一致。
\end{alertblock}
\end{columns}
\end{frame}

% ----------------------------------------------------------------------
\begin{frame}{Week 2 在八周项目中的位置}
\centering
\begin{tikzpicture}[node distance=0.55cm and 0.48cm, every node/.style={font=\scriptsize}]
\tikzset{box/.style={draw, rounded corners, align=center, minimum width=2.12cm, minimum height=0.68cm, fill=blue!6}}
\tikzset{active/.style={draw=darkblue, very thick, rounded corners, align=center, minimum width=2.12cm, minimum height=0.68cm, fill=blue!18}}
\tikzset{arrow/.style={-{Latex[length=2mm]}, thick}}
\node[box] (w1) {Week 1\\balanced pandapower};
\node[active, right=of w1] (w2) {Week 2\\三相不平衡潮流};
\node[box, right=of w2] (w3) {Week 3\\场景生成};
\node[box, right=of w3] (w4) {Week 4\\N-1 标签};
\node[box, below=of w1] (w5) {Week 5\\ML baseline};
\node[box, right=of w5] (w6) {Week 6\\PI-MLP};
\node[box, right=of w6] (w7) {Week 7\\Phase-aware GNN};
\node[box, right=of w7] (w8) {Week 8\\论文和展示};
\draw[arrow] (w1) -- (w2);
\draw[arrow] (w2) -- (w3);
\draw[arrow] (w3) -- (w4);
\draw[arrow] (w4) -- ++(0,-0.48) -| (w5);
\draw[arrow] (w5) -- (w6);
\draw[arrow] (w6) -- (w7);
\draw[arrow] (w7) -- (w8);
\end{tikzpicture}
\vspace{0.75em}
\begin{block}{为什么 Week 2 很重要?}
后续 AI 模型的输入特征、标签和物理 loss 都来自三相潮流结果。如果三相物理量没有理解清楚, 后续模型只是黑箱拟合。
\end{block}
\end{frame}

% ----------------------------------------------------------------------
\begin{frame}{从 Week 1 到 Week 2: 模型复杂度变化}
\begin{tabular}{p{0.30\textwidth}p{0.30\textwidth}p{0.30\textwidth}}
\toprule
项目 & Week 1 & Week 2 \\
\midrule
潮流函数 & \code{pp.runpp()} & \code{pp.runpp\_3ph()} \\
系统假设 & 三相平衡 & 三相不平衡 \\
负荷模型 & \code{load} & \code{asymmetric\_load} \\
DER 模型 & \code{sgen} & \code{asymmetric\_sgen} \\
电压结果 & \code{vm\_pu} & \code{vm\_a/b/c\_pu} \\
线路结果 & 总 loading & A/B/C 相电流与 loading \\
新增指标 & - & VUF, 相别越限 \\
\bottomrule
\end{tabular}
\vspace{0.8em}
\begin{alertblock}{本周必须形成的习惯}
每次得到 \code{res\_*\_3ph} 后, 都要问: 这些结果是否满足对称分量、功率平衡和线路电流公式?
\end{alertblock}
\end{frame}

% ----------------------------------------------------------------------
\section{科研动机}

\begin{frame}{为什么微电网需要三相不平衡分析?}
\begin{columns}[T]
\column{0.53\textwidth}
低压微电网中常见现象:
\begin{itemize}
  \item 单相居民负荷随机接入 A/B/C 相;
  \item 单相 PV, EV charger, heat pump 造成相间差异;
  \item BESS 或 inverter 控制可能改变某一相无功支撑;
  \item 辐射型 feeder 的末端电压更敏感;
  \item islanded mode 下等效电源更弱, 电压更容易波动。
\end{itemize}
\column{0.41\textwidth}
\begin{block}{平衡模型会漏掉什么?}
\begin{itemize}
  \item 单相低电压;
  \item 单相线路过流;
  \item 负序电压过高;
  \item 相别相关的 DER 反向潮流;
  \item 相别相关的 N-1 风险。
\end{itemize}
\end{block}
\end{columns}
\end{frame}

% ----------------------------------------------------------------------
\begin{frame}{本课题的三相安全预测任务}
\centering
\begin{tikzpicture}[node distance=0.80cm and 0.80cm, every node/.style={font=\scriptsize}]
\tikzset{box/.style={draw, rounded corners, align=center, minimum width=2.8cm, minimum height=0.85cm, fill=blue!6}}
\tikzset{arrow/.style={-{Latex[length=2mm]}, thick}}
\node[box] (x) {运行状态 $x_t$\\逐相负荷/PV/BESS};
\node[box, right=of x] (pf) {三相潮流\\\code{runpp\_3ph}};
\node[box, right=of pf] (res) {三相结果\\$V_{i,\phi}, I_{\ell,\phi}$};
\node[box, below=of res] (label) {安全标签\\voltage/current/VUF};
\node[box, left=of label] (ai) {AI screening\\MLP/GNN};
\draw[arrow] (x) -- (pf);
\draw[arrow] (pf) -- (res);
\draw[arrow] (res) -- (label);
\draw[arrow] (x) |- (ai);
\draw[arrow] (ai) -- (label);
\end{tikzpicture}
\vspace{0.8em}
\begin{block}{Week 2 的角色}
本周只完成从 $x_t$ 到三相结果的可信计算。后续 Week 4 再在 contingency $c$ 后重复这一流程生成 $y(x_t,c)$。
\end{block}
\end{frame}

% ----------------------------------------------------------------------
\begin{frame}{三相静态安全指标}
对每个 bus $i$, line $\ell$, phase $\phi\in\{a,b,c\}$:
\vspace{0.2em}
\begin{align*}
\text{Voltage violation:}\quad & |V_{i,\phi}| < V_i^{\min} \quad \text{or} \quad |V_{i,\phi}| > V_i^{\max} \\
\text{Current violation:}\quad & |I_{\ell,\phi}| > I_{\ell}^{\max} \\
\text{Unbalance violation:}\quad & \mathrm{VUF}_{i} > \mathrm{VUF}^{\max}
\end{align*}
\vspace{0.2em}
\begin{block}{教学阈值}
本周 notebook 用 $0.95\le V\le1.05$, line loading $\le100\%$, VUF $\le2\%$ 作为演示阈值。论文中需要说明阈值来源。
\end{block}
\end{frame}

% ----------------------------------------------------------------------
\section{三相物理基础}

\begin{frame}{三相电压: line-to-line 与 line-to-neutral}
\begin{columns}[T]
\column{0.52\textwidth}
低压三相系统常用:
\begin{itemize}
  \item $V_{LL}$: 线电压, 例如 0.4 kV;
  \item $V_{LN}$: 相电压, 约 $0.4/\sqrt{3}$ kV;
  \item pandapower bus 的 \code{vn\_kv} 对 LV feeder 通常表示 nominal line-to-line voltage;
  \item 三相结果表中的 $vm_{a/b/c}$ 是 per-unit 相电压幅值。
\end{itemize}
\column{0.42\textwidth}
\begin{block}{电流复核常用关系}
\[
|I_{\phi}|=\frac{|S_{\phi}|}{|V_{LN,\phi}|}
\]
\[
V_{LN,base}=\frac{V_{LL,base}}{\sqrt{3}}
\]
\end{block}
\end{columns}
\end{frame}

% ----------------------------------------------------------------------
\begin{frame}{相量表示}
每一相电压可以写成复数相量:
\[
V_a = |V_a|e^{j\theta_a},\quad
V_b = |V_b|e^{j\theta_b},\quad
V_c = |V_c|e^{j\theta_c}
\]
\vspace{0.3em}
平衡三相系统中:
\[
|V_a|=|V_b|=|V_c|,\quad
\theta_b\approx\theta_a-120^\circ,\quad
\theta_c\approx\theta_a+120^\circ
\]
\vspace{0.3em}
不平衡系统中, 幅值和相角都可能偏离上述关系。
\begin{alertblock}{AI 特征提示}
只用三相平均电压会丢掉相别风险; phase-aware AI 应保留 A/B/C 通道。
\end{alertblock}
\end{frame}

% ----------------------------------------------------------------------
\begin{frame}{对称分量: 从 ABC 到 012}
定义:
\[
a=e^{j2\pi/3}
\]
\[
\begin{bmatrix}
V^{(0)} \\
V^{(1)} \\
V^{(2)}
\end{bmatrix}
=
\frac{1}{3}
\begin{bmatrix}
1 & 1 & 1 \\
1 & a & a^2 \\
1 & a^2 & a
\end{bmatrix}
\begin{bmatrix}
V_a \\
V_b \\
V_c
\end{bmatrix}
\]
\begin{itemize}
  \item $V^{(1)}$: 正序, 正常三相供电的主要分量;
  \item $V^{(2)}$: 负序, 衡量不平衡程度;
  \item $V^{(0)}$: 零序, 与接地回路和中性线相关。
\end{itemize}
\end{frame}

% ----------------------------------------------------------------------
\begin{frame}{VUF: Voltage Unbalance Factor}
\[
\mathrm{VUF}=100\frac{|V^{(2)}|}{|V^{(1)}|}
\]
\vspace{0.5em}
\begin{columns}[T]
\column{0.48\textwidth}
\begin{block}{为什么需要 VUF?}
\begin{itemize}
  \item 电压幅值都在 0.95-1.05 内, 仍可能存在较强相间不平衡;
  \item inverter, motor, protection settings 都可能受负序分量影响;
  \item VUF 是后续 security label 的重要组成。
\end{itemize}
\end{block}
\column{0.46\textwidth}
\begin{alertblock}{Notebook proof}
从 \code{res\_bus\_3ph} 的 A/B/C 电压幅值和角度手算 VUF, 与 \code{unbalance\_percent} 比较。
\end{alertblock}
\end{columns}
\end{frame}

% ----------------------------------------------------------------------
\begin{frame}{逐相功率和符号约定}
\begin{tabular}{p{0.25\textwidth}p{0.32\textwidth}p{0.34\textwidth}}
\toprule
元素 & pandapower 表 & 本周解释 \\
\midrule
逐相负荷 & \code{asymmetric\_load} & 负荷采用 consumer convention, 正 $P$ 表示消耗 \\
逐相 DER/PV & \code{asymmetric\_sgen} & 正 $P$ 表示发电注入 \\
PCC/slack & \code{ext\_grid} & 结果表中正 $P$ 通常表示外部电网向系统供电 \\
线路损耗 & \code{res\_line\_3ph} & 每相 $pl_\phi$, $ql_\phi$ 可用于功率平衡验证 \\
\bottomrule
\end{tabular}
\vspace{0.8em}
\begin{block}{逐相有功平衡}
\[
P_{grid,\phi}+P_{sgen,\phi}-P_{load,\phi}-P_{loss,\phi}\approx 0
\]
\end{block}
\end{frame}

% ----------------------------------------------------------------------
\section{pandapower 三相建模}

\begin{frame}{\texttt{runpp\_3ph()} 的基本思想}
\begin{columns}[T]
\column{0.50\textwidth}
\code{runpp\_3ph()} 用于 unbalanced/asymmetric/three-phase load flow。
\begin{itemize}
  \item 使用 sequence-frame 建模;
  \item 正序网络使用 Newton-Raphson;
  \item 零序和负序网络使用 current injection 思路;
  \item 三相结果会写入 \code{res\_*\_3ph} 表。
\end{itemize}
\column{0.44\textwidth}
\begin{block}{课程使用方式}
\begin{enumerate}
  \item 建立包含零序参数的 feeder;
  \item 添加逐相负荷和逐相 PV;
  \item 运行 \code{pp.runpp\_3ph(net)};
  \item 对结果做物理 proof。
\end{enumerate}
\end{block}
\end{columns}
\end{frame}

% ----------------------------------------------------------------------
\begin{frame}{线路参数: 为什么需要 zero-sequence?}
三相潮流不能只给正序线路参数。对 line, 本周需要:
\begin{itemize}
  \item 正序/负序: \code{r\_ohm\_per\_km}, \code{x\_ohm\_per\_km}, \code{c\_nf\_per\_km};
  \item 零序: \code{r0\_ohm\_per\_km}, \code{x0\_ohm\_per\_km}, \code{c0\_nf\_per\_km};
  \item 热限值: \code{max\_i\_ka};
  \item 长度: \code{length\_km} 必须大于 0。
\end{itemize}
\begin{alertblock}{常见错误}
只复制 Week 1 的 \code{create\_line\_from\_parameters()} 而不加 zero-sequence 参数, 三相潮流会缺少关键物理信息。
\end{alertblock}
\end{frame}

% ----------------------------------------------------------------------
\begin{frame}[fragile]{三相 feeder 的核心建模代码}
\begin{lstlisting}
net = pp.create_empty_network(sn_mva=1.0)
b0 = pp.create_bus(net, vn_kv=0.4, name="PCC")
b1 = pp.create_bus(net, vn_kv=0.4, name="Bus 1")

pp.create_ext_grid(net, bus=b0, vm_pu=1.0,
                   s_sc_max_mva=10, s_sc_min_mva=8,
                   rx_max=0.1, rx_min=0.1,
                   x0x_max=1.0, x0x_min=1.0,
                   r0x0_max=0.1, r0x0_min=0.1)

pp.create_line_from_parameters(
    net, from_bus=b0, to_bus=b1, length_km=0.08,
    r_ohm_per_km=0.642, x_ohm_per_km=0.083,
    c_nf_per_km=210, max_i_ka=0.20,
    r0_ohm_per_km=1.80, x0_ohm_per_km=0.25,
    c0_nf_per_km=70)
\end{lstlisting}
\end{frame}

% ----------------------------------------------------------------------
\begin{frame}[fragile]{逐相负荷与逐相 PV}
\begin{lstlisting}
pp.create_asymmetric_load(
    net, bus=b1,
    p_a_mw=0.010, p_b_mw=0.008, p_c_mw=0.006,
    q_a_mvar=0.0020, q_b_mvar=0.0016, q_c_mvar=0.0012,
    type="wye", name="unbalanced_load")

pp.create_asymmetric_sgen(
    net, bus=b1,
    p_a_mw=0.003, p_b_mw=0.000, p_c_mw=0.002,
    q_a_mvar=0.0, q_b_mvar=0.0, q_c_mvar=0.0,
    type="wye", name="phase_specific_PV")
\end{lstlisting}
\begin{block}{讲课强调}
负荷和发电不要用“负负荷”混着写。负荷用 \code{asymmetric\_load}, PV/DER 用 \code{asymmetric\_sgen}。
\end{block}
\end{frame}

% ----------------------------------------------------------------------
\begin{frame}{Week 2 教学 feeder}
\centering
\begin{tikzpicture}[node distance=1.0cm and 1.3cm, every node/.style={font=\scriptsize}]
\tikzset{bus/.style={draw, circle, fill=blue!8, minimum size=0.8cm}}
\tikzset{load/.style={draw, rounded corners, fill=orange!12, align=center, minimum width=1.65cm}}
\tikzset{pv/.style={draw, rounded corners, fill=green!12, align=center, minimum width=1.65cm}}
\tikzset{arrow/.style={-{Latex[length=2mm]}, thick}}
\node[bus] (b0) {PCC};
\node[bus, right=of b0] (b1) {Bus 1};
\node[bus, right=of b1] (b2) {Bus 2};
\draw[thick] (b0) -- node[above]{Line 0} (b1);
\draw[thick] (b1) -- node[above]{Line 1} (b2);
\node[load, below=0.75cm of b1] (l1) {A/B/C 负荷\\轻度不平衡};
\node[load, below=0.75cm of b2] (l2) {B 相较重\\末端负荷};
\node[pv, above=0.75cm of b2] (pv) {A/C 相 PV\\逐相注入};
\draw[arrow] (l1) -- (b1);
\draw[arrow] (l2) -- (b2);
\draw[arrow] (pv) -- (b2);
\end{tikzpicture}
\vspace{0.6em}
\begin{block}{为什么网络很小?}
小网络适合逐项 proof。等学生会验证后, Week 3 再切到更大 feeder 和场景生成。
\end{block}
\end{frame}

% ----------------------------------------------------------------------
\begin{frame}[fragile]{运行三相潮流}
\begin{lstlisting}
pp.runpp_3ph(
    net,
    calculate_voltage_angles=True,
    init="auto",
    max_iteration=50,
    tolerance_mva=1e-8,
)
\end{lstlisting}
\vspace{0.5em}
运行后重点查看:
\begin{itemize}
  \item \code{net.converged};
  \item \code{net.res\_bus\_3ph};
  \item \code{net.res\_line\_3ph};
  \item \code{net.res\_ext\_grid\_3ph};
  \item \code{net.res\_asymmetric\_load\_3ph};
  \item \code{net.res\_asymmetric\_sgen\_3ph}。
\end{itemize}
\end{frame}

% ----------------------------------------------------------------------
\begin{frame}{三相结果表怎么读?}
\begin{tabular}{p{0.30\textwidth}p{0.62\textwidth}}
\toprule
结果表 & 课堂重点 \\
\midrule
\code{res\_bus\_3ph} & 每个 bus 的 \code{vm\_a/b/c\_pu}, \code{va\_a/b/c\_degree}, \code{unbalance\_percent} \\
\code{res\_line\_3ph} & 每条 line 的逐相 P/Q, current, loading, loss \\
\code{res\_ext\_grid\_3ph} & PCC 各相注入 P/Q, 用于功率平衡 \\
\code{res\_asymmetric\_load\_3ph} & 逐相负荷实际 P/Q \\
\code{res\_asymmetric\_sgen\_3ph} & 逐相 DER/PV 实际 P/Q \\
\bottomrule
\end{tabular}
\vspace{0.8em}
\begin{alertblock}{不要只看总量}
三相安全预测中, 最危险的经常是某一相, 不是三相平均值。
\end{alertblock}
\end{frame}

% ----------------------------------------------------------------------
\section{Proof 与交叉验证}

\begin{frame}{为什么 notebook 必须做 proof?}
\begin{columns}[T]
\column{0.52\textwidth}
科研代码常见问题:
\begin{itemize}
  \item 负荷和发电符号写反;
  \item kV 和 pu 的基准混淆;
  \item line-to-line 与 line-to-neutral 混淆;
  \item 只检查总功率, 没有逐相检查;
  \item 模型训练前标签就已经错了。
\end{itemize}
\column{0.42\textwidth}
\begin{block}{本周 proof 列表}
\begin{enumerate}
  \item 输入结构检查;
  \item VUF 手算复核;
  \item 逐相 P/Q 功率平衡;
  \item $S=VI^*$ 电流复核;
  \item loading percent 复核;
  \item balanced runpp 交叉验证。
\end{enumerate}
\end{block}
\end{columns}
\end{frame}

% ----------------------------------------------------------------------
\begin{frame}{Proof 0: 输入结构检查}
运行潮流前检查:
\begin{itemize}
  \item 只有一个 PCC/slack;
  \item 所有 element 的 bus index 都存在;
  \item line 参数长度和电流限值为正;
  \item line 包含 zero-sequence 参数;
  \item \code{asymmetric\_load} 和 \code{asymmetric\_sgen} 包含 A/B/C 的 P/Q 列;
  \item \code{ext\_grid} 包含三相潮流需要的短路/序参数字段。
\end{itemize}
\begin{block}{教学目的}
先让输入数据“可被证明为合理”, 再讨论潮流输出。
\end{block}
\end{frame}

% ----------------------------------------------------------------------
\begin{frame}{Proof 1: VUF 手算复核}
从 \code{res\_bus\_3ph} 构造相量:
\[
V_\phi = vm_\phi e^{j\theta_\phi}
\]
再用对称分量矩阵得到:
\[
V^{(1)},\quad V^{(2)}
\]
最后比较:
\[
100\frac{|V^{(2)}|}{|V^{(1)}|}
\quad \text{vs.}\quad
\code{unbalance\_percent}
\]
\begin{alertblock}{断言}
Notebook 中若最大误差超过容差, 自动抛出 \code{AssertionError}。
\end{alertblock}
\end{frame}

% ----------------------------------------------------------------------
\begin{frame}{Proof 2: 逐相功率平衡}
每一相分别验证:
\[
P_{grid,\phi}+P_{sgen,\phi}-P_{load,\phi}-P_{loss,\phi}\approx 0
\]
\[
Q_{grid,\phi}+Q_{sgen,\phi}-Q_{load,\phi}-Q_{loss,\phi}\approx 0
\]
\vspace{0.5em}
\begin{block}{为什么不是只看三相总功率?}
三相总功率平衡并不能证明逐相建模正确。相别交换、某一相符号错误, 都可能被总量掩盖。
\end{block}
\end{frame}

% ----------------------------------------------------------------------
\begin{frame}{Proof 3: 从 $S=VI^*$ 复核相电流}
对线路每一端、每一相:
\[
S_{\phi}=P_{\phi}+jQ_{\phi}=V_{\phi}I_{\phi}^{*}
\]
\[
|I_{\phi}|=\frac{|S_{\phi}|}{|V_{LN,\phi}|}
\]
单位上:
\[
\frac{\mathrm{MVA}}{\mathrm{kV}}=\mathrm{kA}
\]
\begin{alertblock}{关键换算}
\[
V_{LN,base}=\frac{V_{LL,base}}{\sqrt{3}}
\]
\end{alertblock}
\end{frame}

% ----------------------------------------------------------------------
\begin{frame}{Proof 4: 复核 line loading}
每相 loading:
\[
loading_{\ell,\phi}=100\frac{|I_{\ell,\phi}|}{I_\ell^{max}}
\]
线路总 loading:
\[
loading_\ell=\max_{\phi\in\{a,b,c\}}loading_{\ell,\phi}
\]
\vspace{0.5em}
\begin{block}{为什么这很重要?}
后续 thermal violation label 直接来自 loading。如果 loading 的定义没搞清楚, AI 学到的标签就没有物理可信度。
\end{block}
\end{frame}

% ----------------------------------------------------------------------
\begin{frame}{Cross-validation: 平衡 case 的双求解器/双模型复核}
构造一个三相完全平衡的 case:
\begin{itemize}
  \item A/B/C 负荷相同;
  \item A/B/C PV 相同;
  \item 网络参数保持一致。
\end{itemize}
\vspace{0.3em}
然后比较:
\begin{itemize}
  \item \code{runpp\_3ph()} 的三相平均 bus voltage;
  \item \code{runpp()} 的 balanced bus voltage;
  \item \code{runpp\_3ph()} 的 line loading;
  \item \code{runpp()} 的 line loading。
\end{itemize}
\begin{alertblock}{预期结果}
误差应接近数值容差; 三相电压 spread 应接近 0。
\end{alertblock}
\end{frame}

% ----------------------------------------------------------------------
\section{场景实验}

\begin{frame}{Week 2 场景实验}
\begin{tabular}{p{0.24\textwidth}p{0.62\textwidth}}
\toprule
场景 & 目的 \\
\midrule
\code{base} & 基准不平衡场景 \\
\code{high\_load} & 检查负荷升高是否导致电压下降 \\
\code{phase\_b\_heavy\_load} & 检查 B 相重载是否导致 B 相更低电压和更大 VUF \\
\code{single\_phase\_pv\_a} & 检查 A 相 PV 增大是否降低 PCC 进口并改变电压 \\
\code{stress\_unbalance} & 观察强不平衡下的 VUF 和相电压风险 \\
\bottomrule
\end{tabular}
\end{frame}

% ----------------------------------------------------------------------
\begin{frame}{场景 sanity checks}
Notebook 会自动检查:
\begin{enumerate}
  \item \code{high\_load} 的最小电压低于 \code{base};
  \item \code{single\_phase\_pv\_a} 的 PCC 有功进口低于 \code{base};
  \item \code{stress\_unbalance} 的最大 VUF 高于 \code{base};
  \item \code{phase\_b\_heavy\_load} 的 B 相最低电压低于 \code{base}。
\end{enumerate}
\vspace{0.6em}
\begin{block}{为什么叫 sanity check?}
它不能代替严格验证, 但可以快速发现最常见的符号、缩放、相别映射错误。
\end{block}
\end{frame}

% ----------------------------------------------------------------------
\begin{frame}{三相安全 summary: 后续标签的雏形}
每个场景导出:
\begin{itemize}
  \item \code{min\_vm\_a/b/c\_pu};
  \item \code{min\_vm\_all\_pu};
  \item \code{max\_vm\_all\_pu};
  \item \code{max\_unbalance\_percent};
  \item \code{max\_loading\_a/b/c\_percent};
  \item \code{max\_loading\_percent};
  \item \code{p\_grid\_total\_mw};
  \item \code{voltage/loading/unbalance\_violation}。
\end{itemize}
\begin{alertblock}{Week 4 扩展}
把 summary 从 base-case 扩展到每个 $(scenario, contingency)$, 就得到 N-1 label dataset。
\end{alertblock}
\end{frame}

% ----------------------------------------------------------------------
\begin{frame}{从三相潮流到 AI 特征}
\begin{columns}[T]
\column{0.48\textwidth}
\begin{block}{Node features}
\begin{itemize}
  \item $P_{load,a/b/c}, Q_{load,a/b/c}$;
  \item $P_{PV,a/b/c}, Q_{PV,a/b/c}$;
  \item $V_{a/b/c}$, $\theta_{a/b/c}$;
  \item VUF, voltage margin。
\end{itemize}
\end{block}
\column{0.48\textwidth}
\begin{block}{Edge/global features}
\begin{itemize}
  \item line $r,x,r_0,x_0$, length, $I^{max}$;
  \item base-case current/loading by phase;
  \item scenario type;
  \item PCC/grid-connected indicator;
  \item future contingency mask。
\end{itemize}
\end{block}
\end{columns}
\end{frame}

% ----------------------------------------------------------------------
\begin{frame}{课堂 Notebook 结构}
\begin{enumerate}
  \item Setup 和教学阈值;
  \item 构建三相教学 feeder;
  \item 输入网络 proof;
  \item 运行 \code{runpp\_3ph()};
  \item 结果表结构检查;
  \item VUF 手算复核;
  \item 逐相 P/Q 平衡复核;
  \item $S=VI^*$ 电流复核;
  \item line loading 复核;
  \item balanced \code{runpp()} vs \code{runpp\_3ph()} 交叉验证;
  \item 多场景 sanity checks;
  \item 导出 Week 2 summary。
\end{enumerate}
\end{frame}

% ----------------------------------------------------------------------
\begin{frame}{学生本周输出}
\begin{itemize}
  \item 一个可运行的三相不平衡潮流 notebook;
  \item 一张 base-case 三相电压图;
  \item 一张 base-case 三相 line loading 图;
  \item 一张 VUF 图;
  \item 一个 \code{week2\_scenario\_security\_summary.csv};
  \item 一段文字解释: 为什么三相平均量不足以做微电网安全预测;
  \item 完成 2-3 个 exercises。
\end{itemize}
\end{frame}

% ----------------------------------------------------------------------
\begin{frame}{学生练习}
\begin{block}{Exercise 1: 健康 base case}
修改负荷或线路长度, 使 \code{min\_vm\_all\_pu > 0.97} 且 \code{max\_unbalance\_percent < 1.0}。
\end{block}
\begin{block}{Exercise 2: 制造电压越限}
修改 stress scenario, 让某一相电压低于 0.90 pu, 但潮流仍收敛。
\end{block}
\begin{block}{Exercise 3: 单相 PV 思考}
增大 A 相 PV 后, 为什么 B/C 相电压也可能变化?
\end{block}
\end{frame}

% ----------------------------------------------------------------------
\begin{frame}{过渡到 Week 3}
Week 3 会把本周小 feeder 扩展为 scenario generation pipeline:
\begin{itemize}
  \item 加入更多 buses/loads/PV/BESS;
  \item 生成 high PV, evening peak, EV charging, BESS charging/discharging 等场景;
  \item 保存 base-case 三相特征;
  \item 训练前先做数据质量检查;
  \item 为 Week 4 的 N-1 扫描准备候选 contingencies。
\end{itemize}
\vspace{0.6em}
\begin{alertblock}{课程主线}
每周都要保留 proof/cross-validation, 不要让数据生成变成不可解释的脚本。
\end{alertblock}
\end{frame}

% ----------------------------------------------------------------------
\begin{frame}{参考资料}
\small
\begin{itemize}
  \item pandapower documentation, \emph{Asymmetric / Three Phase Power Flow}, \url{https://pandapower.readthedocs.io/en/stable/powerflow/ac_3ph.html}
  \item pandapower documentation, \emph{Asymmetric Load}, \url{https://pandapower.readthedocs.io/en/stable/elements/asymmetric_load.html}
  \item pandapower documentation, \emph{Asymmetric Static Generator}, \url{https://pandapower.readthedocs.io/en/stable/elements/asymmetric_sgen.html}
  \item pandapower documentation, \emph{Line}, \url{https://pandapower.readthedocs.io/en/stable/elements/line.html}
\end{itemize}
\vspace{0.5em}
\begin{block}{建议阅读方式}
先跑通 notebook, 再回头查文档中的符号约定和结果表定义。
\end{block}
\end{frame}

\end{document}
